\section{Supersymmetric Gauge Theories}

In this section we will explain some aspects of supersymmetry and
supersymmetric field theories which are relevant to the work of Witten
and Seiberg. We start by explaining our conventions and then briefly
describe the representations of supersymmetry algebra with and without
central charges. We then discuss the representations of $N=1$
supersymmetry in terms of quantum fields and construct Lagrangians
with $N=1$ and $N=2$ supersymmetry. Most of the material in this
section is by now standard and can be found in \cite{WB,West,Sohnius}. 
Towards the end of this section, we will explicitly calculate the
central charges in $N=2$ theories with and without matter.
\subsection{Conventions}

We start by descrbing our conventions. We use the flat metric
$\eta_{ab} = {\rm diag}~ (1,-1,-1,-1)$. The spinors of the Lorentz
group $SL(2, C)\sim SU(2)_L\times SU(2)_R$ are written with dotted and
undotted components and, under $SL(2,C)$, transform as 
$$
\psi'_\alpha = M_{\alpha}^{~\beta} \psi_\beta\,, \qquad
\bar\psi'_{\dot{\alpha}} = M^{*~\dot{\beta}}_{\dot{\alpha}}
\bar{\psi}_{\dot{\beta}}\,.
$$
Spinor indices are raised or lowered with the $\epsilon$-tensor,
$$
\epsilon^{\alpha \beta} = \epsilon^{\dot{\alpha}\dot{\beta}} =
\left( \begin{array}{lr}
0 & 1 \\
-1 & 0
\end{array} \right) =
(i \sigma_{2})\,.
$$
By definition, this tensor is invariant under a $SL(2,C)$
transformation: $\epsilon^{\alpha \beta} = M_{~\gamma}^{\alpha}
\epsilon^{\gamma \delta}M_{\delta}^{~\beta}$. This can be written as 
$M^T\sigma_2 M=\sigma_2$ which implies $\sigma_2 M=(M^T)^{-1}\sigma_2$
Using this, we can write the transformations of the spinors with
raised indices as
$$
\psi'^\alpha =\psi^\beta (M^{-1})^{~\alpha}_{\beta}\,,\qquad
\bar{\psi}'^{\dot{\alpha}} = \bar{\psi}^{\dot{\beta}}
(M^*)^{-1~\dot{\alpha}}_{\dot{\beta}}\,.
$$
Now, let us define
$$
(\sigma^\mu)_{\alpha \dot{\alpha}} \equiv (1, \vec{\sigma})\,,
$$
then, 
$$
\sigma^\mu P_\mu = \left ( \begin{array}{lr}
P^{0} - P^{3} & -P^{1} + i P^{2} \\
-P^{1} - iP^{2} & P^{0} + P^{3}
\end{array} \right )\,, 
$$
and $\det(\sigma^\mu P_\mu)=P_\mu P^\mu$. We can raise the indices on 
$\sigma^\mu$ using the $\epsilon$-tensor and define $\bar\sigma$ as
$$
(\bar{\sigma}^\mu)^{\dot{\alpha}\alpha}=-(\sigma^\mu)^{\alpha\dot{\alpha}}
=\epsilon^{\dot{\alpha}\dot{\beta}}\epsilon^{\alpha \beta}
(\sigma^\mu)_{\beta \dot{\beta}}\,. 
$$
Numerically, this gives,
$$
(\bar{\sigma}^\mu) = (i\sigma_2) (\sigma^\mu)^T (i\sigma_2)^T =
\sigma_2 (\sigma^\mu)^T\sigma_2 = (1, - \vec{\sigma})\,.
$$
With these conventions, Lorentz transformations are generated by
\begin{eqnarray}
(\sigma^{\mu\nu})_\alpha^{~\beta} & = & \frac{1}{4}
[\sigma^\mu_{\alpha \dot{\beta}} \bar{\sigma}^{\nu\dot{\beta}\beta}- 
(\mu \leftrightarrow \nu)]\,, \nonumber \\
(\bar{\sigma}^{\mu\nu})^{\dot{\alpha}}_{\dot{~\beta}}&=&\frac{1}{4}
[\bar{\sigma}^{\mu\dot{\alpha}\beta}\sigma^\nu_{\beta\dot{\beta}}-
(\mu\leftrightarrow \nu)]\,.\nonumber
\end{eqnarray}
For the scalar product of spinors, we use the following conventions 
\begin{eqnarray}
\psi \chi &= & \psi^{\alpha} \chi_{\alpha} = - \psi_{\alpha}
\chi^{\alpha} = \chi^{\alpha} \psi_{\alpha} = \chi \psi\,,\nonumber\\ 
\bar{\psi} \bar{\chi} &=& \bar{\psi}_{\dot{\alpha}} 
\bar{\chi}^{\dot{\alpha}} =\bar{\chi} \bar{\psi}\,,\nonumber \\
(\psi\chi)^\dagger&=&\bar{\chi}_{\dot{\alpha}}\bar{\psi}^{\dot{\alpha}}
=\bar{\chi}\bar{\psi}=\bar{\psi}\bar{\chi}\,. \nonumber
\end{eqnarray}
We list some more spinor identities
$$
\begin{array}{rl}
\chi\sigma^\mu\bar{\psi} & =-\bar{\psi}\bar{\sigma}^\mu\chi\,, \\
(\chi\sigma^\mu\bar{\psi})^\dagger & =\psi\sigma^\mu\bar{\chi}\,, \\
\chi\sigma^\mu\bar{\sigma}^\nu\psi & =\psi\sigma^\nu\bar{\sigma}^\mu
\chi \,,\\ 
(\chi\sigma^\mu\bar{\sigma}^\nu\psi)^\dagger & =\bar{\psi} 
\bar{\sigma}^\nu \sigma^\mu \bar{\chi}\,.
\end{array}
$$
In the above basis, the Dirac matrices and Dirac and Majorana spinors
are given by
\begin{equation}
\gamma^\mu = \left( \begin{array}{cc}
0 & \sigma^\mu \\ \bar{\sigma}^\mu & 0 \end{array} \right)\,, \qquad
\psi_D=\left(\begin{array}{c}\psi_\alpha\\\bar{\chi}^{\dot\alpha}
\end{array}\right)\,,\qquad
\psi_M=\left(\begin{array}{c}\psi_\alpha\\{\bar{\psi}^{\dot{\alpha}}}
\end{array}\right)\,.
\label{sp-comp}
\end{equation}
As usual, one defines $\gamma_5=-i\gamma^0\gamma^1\gamma^2\gamma^3$.
Consider a massless fermion moving in the $z$-direction. Then, $P^\mu=
E(1,0,0,1)$, and the Dirac equation gives $(\gamma^0-\gamma^3)\psi=0$. 
Since the helicity operator is now $J_3=\frac{i}{2}\gamma^1\gamma^2$,
one gets, $J_3\psi=\frac{i}{2}(\gamma^0)^2 \gamma^1\gamma^2\psi =
\frac{i}{2}\gamma^0\gamma^1\gamma^2\gamma^3\psi=-\frac{1}{2}\gamma_5\psi$.
Hence, 
$$
\gamma_5 =+1 \Rightarrow -ve \;{\rm helicity}\,,
\gamma_5 =-1 \Rightarrow +ve \;{\rm helicity}\,.
$$

\subsection{Supersymmetry Algebra without Central Charges}

In the absence of central charges, the supersymmetry algebra is
written as  
\begin{eqnarray}
&&\{ Q^{I}_{\alpha}, \bar{Q}_{\dot{\alpha} J} \} = 2
\sigma^\mu_{\alpha \dot{\alpha}} P_\mu \delta^{I}_{J}\,, 
\nonumber\\
&&\{Q^I_\alpha , Q^J_\beta\}=0\,,\,\{{\bar Q}_{\dot\alpha I}\,,\, 
{\bar Q}_{\dot\beta J}\}=0\,.
\label{susy-noC}
\end{eqnarray}
Here, $Q$ and $\bar Q$ are the supersymmetry generators and transform
as spin-half operators under the  angular momentum algebra. The
indices $I, J$ run from $1$ to $N$, where $N$ is the total number of
supersymmetries. Moreover, the supersymmetry generators commute with
the momentum operator $P_\mu$ and hence, with $P^2$. Therefore, all
states in a given representation of the algebra have the same mass.  
For a theory to be supersymmetric, it is necessary that its particle
content form a representation of the above algebra. The supersymmetry
algebra can be embedded in the super-Poincar\'{e} algebra and its
representations can be obtained systematically using Wigner's
method. In the following, we will give a brief description of the
representations of supersymmetry algebra. 

\noindent\underline{Massless Irreducible Representations}:
\noindent For massless states, we can always go to a frame where
$P^{\mu} = M (1,0,0,1)$. Then the supersymmetry algebra becomes 
$$
\{ Q^{I}_{\alpha}, \bar{Q}_{\dot{\alpha}J} \} =
\left ( \begin{array}{lr}
0 & 0 \\ 0 & 4M \end{array} \right )
\delta^{I}_{J}\,.
$$
Now, in a unitary theory the norm of a state is always positive
definite. Since $Q_\alpha$ and $\bar Q_{\dot\alpha}$ are conjugate to
each other, and $\{ Q_{1}, \bar{Q}_{\dot{1}}\}=0$, it follows that 
$Q_1|phys>= \bar{Q}_{\dot{1}}|phys> =0$. As for the other generators,
it is convenient to rescale them as  
$$
a^I=\frac{1}{2\sqrt{M}}Q^I_2\,,\qquad
(a^I)^\dagger=\frac{1}{2\sqrt{M}}\bar Q^I_{\dot 2}\,.
$$
Then, the supersymmetry algebra takes the form
$$
\{ a^I, (a^J)^\dagger \} = \delta^{IJ}\,,\quad  \{ a^I, a^J \}=0\,,
\quad \{ (a^I)^\dagger, (a^J)^\dagger \} = 0\,.
$$
This is a Clifford algebra with $2N$ generators and has a
$2^{N}$-dimensional representation. From the point of view of the
angular momentum algebra, $a^I$ is a rising operator and
$(a^I)^\dagger$ is a lowering operator for the helicity of massless
states. We choose the vacuum such that $J_3|\Omega_\lambda>=\lambda 
|\Omega_\lambda>$ and $a^I|\Omega_\lambda>=0$ for all $I$. Other
states are generated by the action of $(a^I)^\dagger$'s on the vacuum
state. From anti-symmetry it follows that a state with $m$
$(a^I)^\dagger$'s, and hence with helicity $\lambda -m/2$, will have a
degeneracy of $^N{\bf C}_m$. The helicity of all states so constructed
will span the range $\lambda$ to $\lambda-N/2$. Some examples are:
$$
\begin{array}{rrrrrr}
N = 1\,:& \quad |\lambda >,&    |\lambda-1/2>\hphantom{,}& & &\\
N = 2\,:& \quad |\lambda >,& 2\,|\lambda-1/2>,&|\lambda -1>\hphantom{,}&&\\ 
N = 4\,:& \quad |\lambda >,& 4\,|\lambda-1/2>,&
6\,|\lambda -1 >,&4\,|\lambda - 3/2>,& |\lambda -2 >
\end{array}
$$
The irreducible representations are not necessarily CPT
invariant. Therefore, if we want to assign physical states to these
representations, we have to suplement them with their CPT conjugates.
If a representation is CPT self-conjugate, it is left unchanged.  
Below, we list the representations after the addition of the CPT
conjugates and indicate the particle spectra which can be assiged to
them:
$$
\begin{array}{lrllllll}
N=1,&\lambda=1/2:&\,\,\, |1/2>\,,&|0>\,,   &|-1/2>\,,&|0>\hphantom{,}&& \\
    &\lambda=1:&\,\,\,|1> \,,&|1/2>\,, &|-1>\,,  &|-1/2>\hphantom{,}&& \\
N=2,&\lambda=1/2:&\,\,\,|1/2>\,,&2|0>\,,&|-1/2>\,,&|-1/2>\,,&2|0>\,,&|1/2>\\
    &\lambda=1:&\,\,\, |1> \,,&2|1/2>\,,&|0>\,,   &|-1>\,, &2|-1/2>\,,&|0>\\
N=4,&\lambda=1:&\,\,\, |1>\,,&4|1/2>\,,&6|0>\,,&4|-1/2>\,,&|-1>\hphantom{,}&
\end{array}
$$
Thus, for $N=1$, the representation contains a Majorana spinor and a
complex scalar if $\lambda=1/2$ (scalar multiplet), or a massless
vector and a Majorana spinor if $\lambda=1$ (vector multiplet). 
For $N=2$ and $\lambda=1/2$, we have two Majorana spinors (or one
Dirac spinor) with two complex scalars. This representation has the
same particle content as two copies of the $N=1$, $\lambda=1/2$
multiplet. For $N=2$ and $\lambda=1$, we have a massless vector, two
Majorana spinors and a complex scalar. Note that this multiplet has
the same particle content as the two $N=1$ multiplets for
$\lambda=1/2$ and $\lambda=1$ put together. For $N=4$, the
representation is self-conjugate and accommodates a massless vector,
two Dirac fermions and three complex scalars.   

\noindent\underline{Massive Irreducible Representations}:
\noindent For massive states, we can always go to the rest frame where
$P_\mu=(M,0,0,0)$ and define 
$$
a^I_\alpha=Q^I_\alpha/\sqrt{2M}\,,\quad (a^I_\alpha)^\dagger= 
{\bar Q}_{\dot\alpha I}/\sqrt{2M}\,.
$$
Then the supersymmetry algebra reduces to 
$$
\{a^I_1,(a_1^J)^\dagger\}=\delta^{IJ}\,,\qquad 
\{a^I_2,(a_2^J)^\dagger\}=\delta^{IJ}\,,
$$
with all other anti-commutators vanishing. The Clifford vacuum is
defined by $a^I_\alpha |\Omega>=0$ and the representation is
constructed by applying $(a^I_\alpha)^\dagger$'s on $\Omega$. Let
$|\Omega>$ be a spin singlet. Then there are $^{2N}{\rm C}_m$ states
at level $m$ and the dimension of the the representation is given by  
$\sum_{m=0}^{2N}{}^{2N}{\rm\bf C}_m=2^{2N}$. The maximum spin which can
be reached is $N/2$ and not $N$ as one might naively expect. This is
because $(a^I_1)^\dagger(a^I_2)^\dagger=\frac{1}{2}
\epsilon^{\alpha\beta} (a^I_\alpha)^\dagger(a^I_\beta)^\dagger$ is a
scalar. Thus the state with $m=2N$ has spin zero, as the vacuum. The
degeneracy of states with a given spin is labelled by the irreducible
representations of the group $USp(2N)$ which we will not discuss
here. Instead, let us consider the simplest example. For $N=1$, the
massive representation contains $2^2=4$ states, 
$$
|\Omega>,\quad a^\dagger_\alpha |\Omega>,\quad\frac{1}{\sqrt{2}}
\epsilon^{\alpha \beta} a^\dagger_\alpha a^\dagger_\beta |\Omega>\,,
$$
with spin content $(0)\oplus (1/2)\oplus (0)$. Here, $(j)$ denotes a
state of total spin $j$ and degeneracy $2j+1$. Thus, in the above
example, we have a Weyl (or Majorana) spinor and a complex scalar
$(\lambda , \phi)$. For $N=2$, the representation contains $2^4=16$
states which, under the $SU(2)$ of angular momentum, decompose as
$5(0)\oplus 4(1/2)\oplus 1(1)$. The $N=4$ massive multiplet has
$2^8=256$ states and inclueds a spin $2$ state.

Till now we have considered representations based on a singlet vacuum.
Let us consider a vacuum $|\Omega_j>$ of spin $j$ which is $2j+1$-fold
degenerate. The representation now contains $(2j+1)2^{2N}$ states. The
spectrum is worked out by combining the $j=0$ representation of the
Clifford algebra and a spin $j$, using the angular momentum addition
rules. For example, to obtain the $N=1$ representation based on
$|\Omega_j>$, we combine $(0)\oplus(1/2)\oplus (0)$ with $(j)$ to
obtain $(j)\oplus(j+1/2)\oplus(j-1/2)\oplus(j)$.  For $j=1/2$, we get
$(1/2)\oplus(1)\oplus(0)\oplus(1/2)$ which corresponds to a gauge
field, a Dirac fermion and a scalar field, all of the same mass.  Note
that in all cases we get the same number of bosonic and fermionic
degrees of freedom.

\subsection{Supersymmetry Algebra with Central Charges}
 
As shown by Haag, Lapuszanski and Sohnius \cite{HLS}, the
supersymmetry algebra (\ref{susy-noC}) admits a central extension and
can be generalised to 
\begin{eqnarray}
&&\{ Q^{I}_{\alpha}, \bar{Q}_{\dot{\beta}J} \} =
2 \sigma^\mu_{\alpha \dot{\beta}} P_\mu \delta^{I}_{J}\,, 
\nonumber\\
&&\{ Q^I_\alpha, Q^J_\beta \}=2{\sqrt 2}\epsilon_{\alpha\beta} Z^{IJ}\,,
\nonumber \\
&&\{\bar{Q}_{\dot{\alpha}I},\bar{Q}_{\dot{\beta}J} \} =
2{\sqrt 2}\epsilon_{\dot{\alpha} \dot{\beta}} Z^{*}_{IJ}\,,
\label{susy-C}
\end{eqnarray}
where, $Z$ and $Z^*$ are the central charge metrices which are
antisymmetric in $I$ and $J$. Let us focus on the case of even $N$.
Using a unitary transformation, we can skew-diagonalize $Z$:
$\tilde{Z}^{IJ} = U^{I}_{A}U^{J}_{B}Z^{AB}$, so that it takes the form
$Z=\epsilon\otimes D$, where $D$ is an $N/2$-dimensional diagonal
matrix. Thus, the index $I$ which counts the number of supersymmetries
can be decomposed into $(a,m)$, with $a=1,2$ coming from the
antisymmetric tensor $\epsilon$, and $m=1,...,N/2$ coming from the
diagonal matrix $D$.  By a further chiral rotation, we may choose the
eigenvalues of $D$ to be real. Once we have skew-diagonalized, it is
sufficient to consider just the $N=2$ supersymmetry, for which the
algebra takes the form
\begin{eqnarray}
&&\{Q^a_\alpha,{\bar Q}_{\dot\beta b}\}=2(\sigma^\mu)_{\alpha\dot\beta}  
P_\mu \delta^{a}_{b}\,,\nonumber \\
&&\{ Q^a_\alpha,Q^b_\beta\}=2{\sqrt 2}\epsilon_{\alpha\beta}
\epsilon^{ab}Z\,, \nonumber \\ 
&&\{{\bar Q}_{\dot\alpha a},{\bar Q}_{\dot\beta b}\}= 
2{\sqrt 2}\epsilon_{\dot\alpha\dot\beta}\epsilon_{ab} Z \,.
\label{Susy-C'}
\end{eqnarray}
Since $Z$ commutes with all the generators, we can fix it to be
the eigenvalue for the given representation. Now, let us define:
$$
a_\alpha=\frac{1}{2} \{ Q^1_\alpha +
\epsilon_{\alpha\beta} (Q^2_\beta)^\dagger \}\,, \qquad
b_\alpha=\frac{1}{2} \{ Q_\alpha^1 -
\epsilon_{\alpha\beta} (Q^2_\beta)^\dagger \}\,. 
$$
Then, the algebra (\ref{susy-C}) reduces to 
\begin{equation}
\{a_\alpha,a^\dagger_\beta\}=\delta_{\alpha\beta}(M+{\sqrt 2}Z)\,,\qquad
\{ b_\alpha,b^\dagger_\beta\}=\delta_{\alpha\beta} (M-{\sqrt 2}Z)\,,
\label{susy-ab}
\end{equation}
with all other anticommutators vanishing. Since all physical states
have positive definite norm, it follows that for massless states, the
central charge is trivially realised ({\it i.e.},\,$Z=0$). For massive
states, this leads to a bound on the mass $M\geq {\sqrt 2}|Z|$. When
$M={\sqrt 2}|Z|$, one set of operators in (\ref{susy-ab}) is trivially
realized and the algebra resembles the massless case and the dimension
of representation is greatly reduced. For example, a reduced massive
$N = 2$ multiplet has the same number of states as a massless $N = 2$
multiplet. Thus the representations of the $N=2$ algebra with a
central charge can be classified as either long multiplets (when $M >
{\sqrt 2}|Z|$) or short multiplets (when $M={\sqrt 2}|Z|$).
 
The mass bound $M\geq {\sqrt 2}|Z|$ is reminiscent of the Bogomol'nyi
bound in the Georgi-Glashow model. In fact, it turns out that in the
supersymmetric version of the Georgi-Glashow model (which is based on
the algebra without central charges) the solitonic solutions do give
rise to a central extension term in the supersymmetry algebra, thus
realizing (\ref{susy-C})\cite{Olive-Witten}. The origin of the central
charge is easy to understand: The supersymmetry charges $Q$ and $\bar
Q$ are space integrals of local expressions in the fields (the time
component of the super-currents). In calculating their anticommutators,
one encounters surface terms which are normally neglected. However, in
the presence of electric and magnetic charges, these surface terms are
non-zero and give rise to a central charge. As we will explicitly show
towards the end of this section, it is found that 
\begin{equation}
Z = a (q+ig)= ae (n+m\tau)\,,
\label{Ztau}
\end{equation}
so that $M\geq {\sqrt 2}|Z|$ coincides with the Bogomol'nyi bound
(\ref{Bbound}). From (\ref{susy-ab}) it is clear that the BPS states 
(which saturate the bound) are annihilated by half of the
supersymmetry generators and thus belong to reduced representations of
(\ref{susy-C}). An important consequence of this is that, for BPS
states, the relationship between their charges and masses is dictated
by supersymmetry and does not receive perturbative or non-perturbative
corrections in quantum theory. This is so because a modification of
this relation implies that the states no longer belong to a short
multiplet. On the other hand, quantum correction are not expected to
generate the extra degrees of freedom needed to convert a short
multiplet into a long multiplet. Since there is no other possibility,
we conclude that for short multiplets the relation $M={\sqrt 2}|Z|$ is
not modified either perturbatively or non-perturbatively.    

\subsection{Local Representations of N=1 Supersymmetry}

In this subsection we describe the action of supersymmetry on the
local fields in a quantum field theory. It is well known that the 
Poincar\'{e} group naturally acts on the space-time coordinates. All
other objects transform as components of tensors or spinors defined on
the space-time manifold. Similarly, the supersymmetry transformations
naturally act on an extension of the space-time, called the
``superspace''. The quantum fields then transform as components of a  
``superfield'' defined on the superspace. In the following, we first
describe these notions and then introduce the chiral and vector
superfields.  

\noindent\underline{Superspace :} 
\noindent The superspace is obtained by adding four spinor degrees of
freedom $\theta^\alpha,\bar\theta_{\dot\alpha}$ to the space-time
coordinates $x^\mu$. The spinor index is raised and lowered with the 
$\epsilon$-tensor and $\theta\theta = \theta^\alpha\theta_\alpha = -2
\theta^1\theta^2$. Similarly, $\bar\theta\bar\theta=
\bar\theta_{\dot\alpha}\bar\theta^{\dot\alpha}= 2\bar\theta_{\dot 1} 
\bar\theta_{\dot 2}$. We also have
$$
\theta^{\alpha}\theta^{\beta}=\frac{1}{2} \epsilon^{\alpha \beta}
\theta \theta\,,\qquad
\bar{\theta}_{\dot{\alpha}}\bar{\theta}_{\dot{\beta}}=-\frac{1}{2}
\epsilon_{\dot{\alpha}\dot{\beta}} \bar{\theta}\bar{\theta}\,,\qquad
\theta \sigma^\mu \bar{\theta} \theta \sigma^\nu \bar{\theta}=
\frac{1}{2}\theta \theta \bar{\theta}\bar{\theta} \eta^{\mu\nu}\,.
$$
These formulae are the basis for Fierz rearrangements. 
 
Under the supersymmetry transformations (\ref{susy-noC}) with $N=1$
and transformation parameters $\xi$ and $\bar\xi$, the superspace
coordinates are taken to transform as  
\begin{eqnarray}
x^\mu &\rightarrow& x'^\mu=x^\mu + i\theta\sigma^\mu\bar\xi
- i\xi\sigma^\mu\bar\theta \,,\nonumber\\
\theta &\rightarrow&\theta'=\theta + \xi\,, \nonumber\\
\bar\theta &\rightarrow& \bar\theta'=\bar\theta +\bar\xi\,.
\label{superspace}
\end{eqnarray}
Since these transformations are implemented by the operator 
$\xi^\alpha Q_\alpha + \bar{\xi}_{\dot\alpha}\bar{Q}^{\dot\alpha}$, we
can easily obtain the representation of the supercharges acting on the
superspace as 
\begin{equation}
Q_\alpha=\frac{\partial}{\partial\theta^\alpha}- 
i\sigma^\mu_{\alpha\dot{\alpha}}\bar{\theta}^{\dot{\alpha}}\,
\partial_\mu\,,\qquad \bar{Q}_{\dot\alpha} = -\frac{\partial}{\partial
\bar{\theta}^{\dot\alpha}}+i\theta^\alpha 
\sigma^\mu_{\alpha\dot{\alpha}}\, \partial_\mu\,.
\label{superspace-Q}
\end{equation}
These satisfy $\{Q_\alpha ,\bar{Q}_{\dot\alpha}\}= 
2i\sigma^\mu_{\alpha\dot\alpha}\,\partial_\mu$.
Moreover, using the chain rule, it is easy to see that 
$\partial/\partial x^\mu$ is invariant under (\ref{superspace}) but
not $\partial/\partial\theta$ and $\partial/\partial\bar\theta$.   
Therefore, we introduce the super-covariant derivatives
\begin{equation}
D_\alpha = \frac{\partial}{\partial \theta^\alpha} +
i\sigma^\mu_{\alpha\dot\alpha}\,\partial_\mu \,,\qquad
\bar{D}_{\dot{\alpha}} = -\frac{\partial}{\partial
\bar{\theta}^{\dot{\alpha}}}-i\sigma^\mu_{\alpha \dot{\alpha}}
\theta^{\alpha} \, \partial_\mu \,.
\label{superspaceD}
\end{equation}
They satisfy $\{ D_\alpha,\bar{D}_{\dot\alpha}\}=-2i\sigma^\mu_{\alpha
\dot{\alpha}}\,\partial_\mu$ and commute with $Q$ and $\bar Q$.

\noindent\underline{Superfields:} 
\noindent A superfield is a function on the superspace, say,
$F(x,\theta,\bar\theta)$. Since the $\theta$-coordinates are
anti-commuting, the most general $N=1$ superfield can always be
expanded as  
\begin{eqnarray}
F(x,\theta,\bar\theta)&=& f(x)+\theta\phi(x) + \bar\theta\bar\chi(x)
+\theta\theta m(x) +\bar\theta\bar\theta n(x) + 
\theta\sigma^\mu\bar\theta v_\mu (x) \nonumber\\
&+&\theta\theta\bar\theta\bar\lambda(x)+\bar\theta\bar\theta\theta\psi(x)
+ \theta\theta\bar\theta\bar\theta d(x)\,. 
\label{superfield}
\end{eqnarray}
Clearly, any function of superfields is, by itself, a superfield. 
Under supersymmetry, the superfield transforms as $\delta F=(\xi
Q+\bar\xi\bar Q)F$, from which, the transformation of the component
fields can be obtained. Note that since $d(x)$ is the component of
highest dimension in the multiplet, its variation under supersymmetry
is always a total derivative of other components. Thus, ignoring
surface terms, the space-time integral of this component is invariant
under supersymmetry. This tells us that a supersymmetric Lagrangian
density may be constructed as the highest dimension component of an
appropriate superfield. To describe physical systems, we do not need
all components of the superfield. The relevant components are selected
by imposing appropriate constraints on the superfield. 

\noindent\underline{Chiral Multiplets:} 
\noindent The $N=1$ scalar multiplet is represented by a superfield
with one constraint:
$$
\bar{D}_{\dot{\alpha}} \Phi = 0 \,.
$$
This is referred to as the chiral superfield. Note that for  
$y^\mu = x^\mu +i\theta\sigma^\mu\bar{\theta}$, we have 
$$
\bar{D}_{\dot\alpha}y^\mu=0,\quad\bar{D}_{\dot\alpha}\theta^\beta=0\,.
$$
Therefore, any function of $(y,\theta)$ is a chiral superfield. It can
be shown that this also is a necessary condition. Hence, any chiral
superfield can be expanded as  
\begin{equation}
\Phi(y,\theta)=A(y)+\sqrt{2}\theta\psi(y)+\theta\theta F(y)\,.
\label{csfy}
\end{equation}
Here, $A$ and $\psi$ are the fermionic and scalar components
respectively and $F$ is an auxiliary field required for the off-shell
closure of the algebra. Similarly, an anti-chiral superfield is
defined by $D_{\alpha}\Phi^\dagger=0$ and can be expanded as 
\begin{equation}
\Phi^\dagger(y^\dagger,\bar\theta) =A^\dagger(y^\dagger)+\sqrt{2}
\bar{\theta}\bar{\psi}(y^\dagger)+\bar{\theta}\bar{\theta} 
F^\dagger(y^\dagger)\,,
\label{acsfy}
\end{equation}
where, $y^{\mu\dagger}=x^\mu-i\theta\sigma^\mu\bar{\theta}$.
The product of chiral superfields is a chiral superfield. In general,
any arbitrary function of chiral superfields is a chiral superfield:
\begin{eqnarray}
{\cal W}(\Phi_{i}) &=& {\cal W}(A_{i} + \sqrt{2} \theta \psi_{i} +
\theta \theta F_{i} ) \nonumber \\
&=&{\cal W}(A_{i}) + \frac{\partial {\cal W}}{\partial A_{i}}
\sqrt{2} \theta \psi_{i} + \theta \theta \left (\frac{\partial 
{\cal W}}{\partial A_{i}} F_{i} - \frac{1}{2} \, \frac{\partial^{2}
{\cal W}}{\partial A_i A_j}\, \psi_{i} \psi_{j} \right)\,.
\label{spotential}
\end{eqnarray}
${\cal W}$ is referred to as the superpotential. In terms of the
original variables, $\Phi$ and $\Phi^\dagger$ take the form 
\begin{eqnarray}
\Phi(x,\theta,\bar\theta) &=& A(x)+i\theta\sigma^\mu \bar{\theta}
\partial_\mu A-\frac{1}{4}\theta^2 \bar{\theta}^2 \Box A 
\nonumber\\
&&+\sqrt{2}\theta\psi(x)-\frac{i}{\sqrt{2}}\theta\theta
\partial_\mu\psi\sigma^\mu\bar{\theta} +\theta \theta F(x)\,,
\label{csfield}
\\
\Phi^\dagger(x,\theta,\bar\theta)&=&A^\dagger(x)-i\theta\sigma^\mu
\bar{\theta} \partial_\mu A^\dagger - \frac{1}{4} \theta^{2}
\bar{\theta}^{2} \Box A^\dagger \nonumber\\
&&+\sqrt{2} \bar{\theta} \bar{\psi} (x) + \frac{i}{\sqrt{2}} \,
\bar{\theta} \bar{\theta} \, \theta \sigma^\mu \, \partial_\mu
\bar{\psi} + \bar{\theta} \bar{\theta} F^\dagger (x)\,.
\label{asfield}
\end{eqnarray}

\noindent\underline{Vector Multiplet:} 
\noindent This multiplet is represented by a real superfield satifying
$V=V^\dagger$. In components, it takes the form
$$
\begin{array}{lll}
V(x, \theta, \bar{\theta}) =& C + i \theta  \chi - i \bar{\theta}
\bar{\chi} + \frac{i}{2} \theta^{2} (M + i N) - \frac{i}{2}
\bar{\theta}^{2} (M - iN) \\
& - \theta \sigma^\mu \bar{\theta} A_\mu + i \theta^{2} \bar{\theta}
(\bar{\lambda} + \frac{i}{2} \bar{\sigma}^\mu \partial_\mu \chi )\\
& - i \bar{\theta}^{2} \theta (\lambda + \frac{i}{2} \sigma^\mu
\partial_\mu \bar{\chi} ) + \frac{1}{2} \theta^{2} \bar{\theta}^{2}
(D - \frac{1}{2} \Box C)\,.
\end{array}
$$
Many of these components can be gauged away using the abelian gauge
transformation $V \rightarrow V + \Lambda + \Lambda^\dagger$, where  
$\Lambda$ ($\Lambda^\dagger$) are chiral (antichiral) superfields.  
In the so called Wess-Zumino gauge, we set $C=M=N=\chi= 0$, so that
$$
V=-\theta \sigma^\mu\bar{\theta}A_\mu+i\theta^{2}
\bar{\theta}\bar{\lambda} - i \bar{\theta}^{2} \theta \lambda +
\frac{1}{2} \theta^{2} \bar{\theta}^{2} D\,.
$$
In this gauge, $V^2=\frac{1}{2}A_\mu A^\mu \theta^{2}\bar{\theta}^{2}$
and $V^{3}=0$. The Wess-Zumino gauge breaks supersymmetry, but not the
gauge symmetry of the abelian gauge field $A_\mu$. The Abelian field
strength is defined by  
$$
W_{\alpha} = - \frac{1}{4} \bar{D}^{2} D_{\alpha} V\,,\qquad
\bar{W}_{\dot{\alpha}}=-\frac{1}{4}D^{2} \bar{D}_{\dot{\alpha}}V\,.
$$
$W_{\alpha}$ is a chiral superfield. Since it is gauge invariant, it
can be computed in the Wess-Zumino gauge and takes the form
\begin{equation}
W_\alpha =-i\lambda_\alpha (y)+\theta_\alpha D-\frac{i}{2} 
(\sigma^\mu\bar{\sigma}^\nu \theta)_\alpha \, F_{\mu\nu}+ 
\theta^2 (\sigma^\mu \partial_\mu \bar{\lambda} )_{\alpha}\,,
\label{W-abelian}
\end{equation}
where, $F_{\mu\nu}=\partial_\mu A_\nu -\partial_\nu A_\mu$ is the
familiar abelian field strength tensor. 

In the non-Abelian case, $V$ belongs to the adjoint representation of
the gauge group: $V=V_AT^A$, where, $T^{A\dagger}=T^A$. The gauge
transformations are now implemented by 
$$
e^{-2V} \rightarrow e^{-i \Lambda^\dagger} e^{-2V} e^{i \Lambda}
\qquad {\rm where}, \; \Lambda = \Lambda_{A}T^{A}
$$
The non-Abelian gauge field strength is defined by 
$$
W_{\alpha} = \frac{1}{8} \bar{D}^{2} e^{2V} D_{\alpha} e^{-2V}
$$
and transforms as 
$$
W_\alpha\rightarrow W'_\alpha= e^{-i\Lambda}W_\alpha e^{i\Lambda}\,.
$$
In components, it takes the form
\begin{equation}
W_\alpha =T^a \left (-i \lambda^{a}_{\alpha} + \theta_{\alpha}
D^{a} - \frac{i}{2}(\sigma^\mu \bar{\sigma}^\nu \theta )_{\alpha}
F^{a}_{\mu\nu}+\theta^2 \sigma^\mu D_\mu\bar{\lambda}^a \right)\,, 
\label{W-nonabelian}
\end{equation}
where, 
$$
F^a_{\mu\nu}=\partial_\mu A^a_\nu -\partial_\nu A^a_\mu + 
\, f^{abc} A^b_\mu A^c_\nu\,,\qquad
D_\mu\bar{\lambda}^a=\partial_\mu\bar{\lambda}^a +
\,f^{abc} A_\mu^b \bar{\lambda}^c \,.
$$
In the next section, we will construct supersymmetric Lagrangians in
terms of superfields. 

\subsection{Construction of N=1 Lagrangians}

In this section we will construct the $N=1$ Lagrangians for the 
scalar and the vector multiplets. These serve as the building blocks
for the $N=2$ Lagrangian which is our real interest. As stated before,
a supersymmetric Lagrangian can be constructed as the highest
component of a superfield. Thus the problem reduces to that of finding
appropriate superfields. 

\noindent\underline{Lagrangian for the Scalar Multiplet:}
\noindent Let us first consider the product of a chiral and an
anti-chiral superfield $\Phi_i^\dagger\Phi_j$. This is a general
superfield and its highest component can be computed using
(\ref{csfield}) and  (\ref{asfield}) as  
\begin{eqnarray}
\Phi^\dagger_i\Phi_j\mid_{\theta^{2}\bar{\theta}^{2}}\,\,\,=
&-&\frac{1}{4}\,A^\dagger_i\Box A_j -\frac{1}{4}\Box A^\dagger_i 
A_j+F^\dagger_i F_j +\frac{1}{2}\,\partial_\mu A^\dagger_i\partial^\mu
A_j \nonumber\\
&-& \frac{i}{2}\,\psi_j\sigma^\mu \partial_\mu \bar{\psi}_{i} +
\frac{i}{2} \partial_\mu \psi_j \sigma^\mu\bar{\psi}_i\,.\nonumber
\end{eqnarray}
Dropping some total derivatives and summing over $i = j$, we get the
free field Lagrangian
$$
{\cal L}=
\Phi^\dagger_i \Phi_i \mid_{\theta^2\bar\theta^2}
=\partial_\mu A^\dagger_i \partial^\mu A_i + F^\dagger_i F_i -i
\bar{\psi}_i \bar{\sigma}^\mu \partial_\mu \psi_i\,.
$$
This is the free Lagrangian for a massless scalar and a massless
fermion with an auxiliary field which can be eliminated by its
equation of motion. Supersymmetric interaction terms can be
constructed in terms of the superpotential (\ref{spotential}) and its
conjugate, which are holomorphic functions of $\Phi$ and
$\Phi^\dagger$, respectively. Moreover, note that the space of the
fields $\Phi$ may have a non-trivial metric $g^{ij}$ in which case the scalar
kinetic term, for example, takes the form $g^{ij}\partial_\mu
A_i^\dagger \partial^\mu A_j$, with appropriate modifications for
other terms. In such cases, the free
field Lagrangian above has to be replaced by a non-linear
$\sigma$-model. Thus, the most general $N=1$ supersymmetric Lagrangian
for the scalar multiplet (including the interaction terms) is given by 
$$
{\cal L}=\int\,d^{4}\theta\,K(\Phi,\Phi^\dagger) +
\int d^{2} \theta {\cal W}(\Phi) + \int d^{2} \bar{\theta} 
\bar{\cal W}(\Phi^\dagger)\,.
$$
Note that the $\theta$-integrals pick up the highest component of the
superfield and in our conventions, $\int d^2 \theta \theta^2=1$ and 
$\int d^2 \bar\theta \bar\theta^2=1$. In terms of the non-holomorphic
function $K(A, A^\dagger)$, the metric on the field space is given by 
$g^{ij}=\partial^2K/\partial A_i\partial A^\dagger_j$. For this
reason, the function $K(\Phi,\Phi^\dagger)$ is referred to as the
K\"{a}hler potential.    

For a renormalizable theory, the forms of $K$ and ${\cal W}$ are not
arbitrary and are constrained by R-symmetry. This symmetry acts on the
chiral superfields as follows
\begin{eqnarray}
R\,\Phi (x,\theta) = \Phi'(x,\theta) &=& e^{2 in\alpha}
\Phi (x, e^{-i\alpha}\theta)\,, \nonumber \\
R\,\Phi^\dagger(x,\bar\theta ) = \Phi'^\dagger(x,\theta) &=& e^{-2
in\alpha}\Phi^\dagger(x, e^{i \alpha}\bar{\theta})\,. \nonumber
\end{eqnarray}
Under this, the component fields transform as 
$$
\begin{array}{lll}
A & \rightarrow e^{2i n \alpha} A \,,\\
\psi & \rightarrow e^{2i(n-1/2) \alpha} \psi\,, \\
F & \rightarrow e^{2i(n-1) \alpha} F\,.
\end{array}
$$
We refer to $n$ as the R-character. Since $\theta \rightarrow
e^{+i\alpha}\theta$, or $d^2\theta \rightarrow e^{-2 i\alpha}d^2
\theta$, The R-character of the superfields in each term of ${\cal W}$
must add up to one. Similarly, $K$ should be $R$-neutral. The vector
multiplet is real and it has no natural $R$-symmetry. This symmetry
plays an important role in the study of supersymmetric gauge theories
and we will come back to it in the next section.

\noindent\underline{Lagrangians for the Vector Multiplet:}
\noindent As mentioned in the previous section, the Abelian field
strength $W$, given by (\ref{W-abelian}), is a chiral
superfield. Using the expansion there, one can easily compute that 
$$
W^\alpha W_\alpha\mid_{\theta\theta}=-2i\lambda\sigma^\mu\partial_\mu
\bar{\lambda} + D^2 - \frac{1}{2} F^{\mu\nu} F_{\mu\nu} +\frac{i}{4} \,
\epsilon^{\mu\nu\rho\sigma} F_{\mu\nu}F_{\rho\sigma}\,.
$$
Hence, the usual abelian supersymmetric Lagrangian (which does not
contain the $F\widetilde F$ term) is given by    
$$
{\cal L}=\frac{1}{4g^{2}}\left(\int d^2\theta\, W^\alpha W_\alpha +
\int d^2\bar\theta\, \bar{W}_{\dot\alpha} \bar{W}^{\dot\alpha}\right)\,.
$$
Similarly, in the non-Abelian case, using the normalization 
${\rm Tr}T^aT^b=\delta^{ab}$, we have
\begin{equation}
{\rm Tr}(W^\alpha W_\alpha \mid_{\theta\theta})=-2i\lambda^a\sigma^\mu 
D_\mu \bar{\lambda}^a +D^a D^a -\frac{1}{2} F^{a\mu\nu}
F^a_{\mu\nu} +\frac{i}{4}\epsilon^{\mu\nu\rho\sigma} F^a_{\mu\nu} 
F^a_{\rho\sigma}\,,
\label{WW-nonabelian}
\end{equation}
and, hence, the usual non-Abelian supersymmetric Lagrangian (without
the $F\widetilde F$-term) is given by 
$$
{\cal L}=\frac{1}{4g^2}{\rm Tr}\left(\int d^2\theta\, W^\alpha W_\alpha 
+\int d^2\bar\theta\, \bar{W}_{\dot\alpha}\bar{W}^{\dot\alpha}\right)\,.
$$
However, we are interesed in the supersymmetric analogue of the
Lagrangian (\ref{LXXVIII}) which also contains a $\theta$-term.
From (\ref{WW-nonabelian}), it is obvious that the super Yang-Mills
Lagrangian with a $\theta$-term can be written as 
\begin{eqnarray}
{\cal L}&=&\frac{1}{8\pi}{\rm Im}\left(\,\tau\,{\rm Tr}\int d^2\theta\,
W^\alpha W_\alpha \right) \nonumber\\
&=&-\frac{1}{4g^2} F^a_{\mu\nu}F^{a\mu\nu}+ \frac{\theta}{32\pi^2}  
F^a_{\mu\nu}\tilde F^{a\mu\nu}+ \frac{1}{g^2}(\frac{1}{2}D^a D^a 
-i\lambda^a\sigma^\mu D_\mu\bar{\lambda}^a) \,,
\label{tauWW}
\end{eqnarray}
where, $\tau = \theta/2\pi + 4\pi i/g^2$. Note that $\tau$ can be
regarded as a constant chiral superfield.  

\noindent\underline{Interaction Terms and the General $N=1$
  Lagrangian:}
Let the chiral superfields $\Phi_i$ belong to a given
representation of the gauge group in which the generators are
the matrices $T^a_{ij}$. The kinetic energy term
$\Phi_i^\dagger\Phi_i$ is invariant under global gauge transformations
$\Phi^\prime=e^{-i\Lambda} \Phi$. In the local case, to insure that
$\Phi^\prime$ remains a chiral superfield, $\Lambda$ has to be a
chiral superfield. The supersymmetric gauge invariant kinetic energy
term is then given by  $\Phi^\dagger e^{-2V}\Phi$. We are now in a
position to write down the full N=1 supersymmetric Lagrangian as 
\begin{equation}
{\cal L}=\frac{1}{8\pi}\,{\rm Im}\left(\tau{\rm Tr}\int d^\theta
\,W^\alpha W_\alpha \right)+\int d^2\theta d^2\bar\theta\,\Phi^\dagger 
e^{-2V}\Phi+ \int d^2\theta\,{\cal W}+\int d^2\bar\theta\,\bar{\cal W}\,.     
\label{fullSF}
\end{equation}
Note that since each term is separately invariant, the relative
normalisation between the scalar part and the Yang-Mills part is
not fixed by $N=1$ supersymmetry. In the above, we have set the
normalization of the scalar part to one, but later, we will change
this by rescaling the scalar multiplet $\Phi$. In terms of component
fields, the above Lagrangian takes the form  
\begin{eqnarray}
{\cal L}\,=&-&\frac{1}{4g^2}F^a_{\mu\nu}F^{a\mu\nu}+\frac{\theta} 
{32\pi^2}F^a_{\mu\nu}\widetilde F^{a\mu\nu}
-\frac{i}{g^2} \,\lambda^a \sigma^\mu D_\mu \bar{\lambda}^a 
+\frac{1}{2g^2}D^a D^a  \nonumber\\
&+& (\partial_\mu A-iA^a_\mu T^a A)^\dagger(\partial^\mu A-iA^{a\mu} 
T^a A) -i\,\bar{\psi} 
\bar{\sigma}^\mu (\partial_\mu\psi -iA^a_\mu T^a \psi)\nonumber\\
&-& D^a A^\dagger T^a A -i\sqrt{2}\,A^\dagger T^a \lambda^a
\psi+i\sqrt{2}\, \bar{\psi} T^a A \bar{\lambda}^a + F^\dagger_i F_i
\nonumber\\ 
&+&\frac{\partial {\cal W}}{\partial A_i} \, F_{i} +
\frac{\partial \bar{\cal W}}{\partial A^\dagger_i} \, F^\dagger_{i} -
\frac{1}{2}\,\frac{\partial^2 {\cal W}}{\partial A_i\partial A_j}\,
\psi_{i}\psi_{j}-\frac{1}{2}\,\frac{\partial^2 \bar{\cal W}} 
{\partial A_i^\dagger\partial
  A_j^\dagger}\,\bar{\psi}_{i}\bar{\psi}_{j}\,. 
\label{fullC}
\end{eqnarray}
Here, ${\cal W}$ denotes the scalar component of the superpotential.
The auxiliary fields $F$ and $D^a$ can be eliminated by using their
equations of motion. The terms involving these fields, thus, give rise 
to the scalar potential  
\begin{equation}
V = \sum_i\,\left|\frac{\partial{\cal W}}{\partial A_i}\right|^2
-\frac{1}{2} g^2 (A^\dagger T^a A)^2\,.
\label{Spot}
\end{equation}
